\documentclass[10pt]{article}

\usepackage[utf8]{inputenc}

\usepackage[T1]{fontenc}

\usepackage{amsmath}

\usepackage{amsfonts}

\usepackage{amssymb}

\usepackage[version=4]{mhchem}

\usepackage{stmaryrd}

\usepackage{bbold}



\begin{document}

замкнутой выпуклой оболочкой $R(U)$ множества $U$ точек



\[

x_{k}=u_{k}(t), \quad k=1,2, \ldots, n, \quad a \leqslant t \leqslant b

\]



(т е о р е м а Р и с с a).



\begin{enumerate}

  \setcounter{enumi}{1}

  \item Каждая точка $x=\left(x_{1}, x_{2}, \ldots, x_{n}\right) \in R(U) \subset \mathbb{R}^{n}$ может быть представлена в виде

\end{enumerate}



\[

x_{k}=\sum_{j=1}^{m} \lambda_{j} u_{k}\left(t_{j}\right), \quad k=1,2, \ldots, n,

\]



где $\lambda_{j}>0, j=1,2, \ldots . m, \sum_{j=1}^{m} \lambda_{j}=1, m \leqslant n+1$, а если $x \in \partial R(U)$, то $m \leqslant n$ (т е о р ем а $\mathbf{K}$ а р а т е о д о р и).



\begin{enumerate}

  \setcounter{enumi}{2}

  \item Для того чтобы существовала, по крайней мере, одна неубывающая функция $\mu(t), a \leqslant t \leqslant b$, такая, что

\end{enumerate}



\[

\int_{a}^{b} w_{k}(t) d \mu(t)=\gamma_{k}, \quad k=1,2, \ldots, n

\]



где



$w_{1}(t) \equiv 1, w_{k}(t)=u_{k}(t)+i v_{k}(t) ; u_{k}(t), v_{k}(t), k=1,2, \ldots, n$, - заданные действительные непрерывные на $\left[\begin{array}{ll}a, & b\end{array}\right]$ функции, $\gamma_{1}>0, \gamma_{k}-$ заданные комплексные числа, необходимо и достаточно, чтобы всякий раз, когда при нек-рых комплексных числах $\alpha_{1}, \alpha_{2}, \ldots, \alpha_{n}$ выполняется неравенство



\[

\sum_{k=1}^{n}\left[\alpha_{k} w_{k}(t)+\bar{\alpha}_{k} \bar{w}_{k}(t)\right] \geqslant 0, \quad a \leqslant t \leqslant b,

\]



имело место также и неравенство



\[

\sum_{k=1}^{n}\left[\alpha_{k} \gamma_{k}+\bar{\alpha}_{k} \bar{\gamma}_{k}\right] \geqslant 0

\]



(те о рем а $\mathrm{P}$ и с с а).



Приведенные теоремы позволили дать геометрич. и алгебраич. характеристики областей значений систем коэффициентов и отдельных коэффициентов на класcах функций, регулярных и имеющих положительную действительную часть в круге (кольце), регулярных и типично вещественных в круге (кольце) и на нек-рых других классах (см. [1] Добавление; [4], [5]). Лит.: [1] $\Gamma$ о лу 3 и н $\Gamma$. М., Геометрическая теория функ-

ций комплексного переменного, 2 изд., М., 1966; [2] К р е й н M. Г., "Успехи матем. наук", 1951, т. 6, в. 4, с. 3-120; [3] Л е б е д е в Н. А., А л е к с а н др о в И. А., "Тр. Матем. там же, с. 33-46; [5] е е Ж е, «Зап. науч. семинаров Јенингр. отделения Матем. ин-та АН СССР», 1972, т. 24, с. $29-62 ; 1974$, т. 44 , с. $17-40 ; 1980$, т. 100 , с. $17-25$.



ПАРАМЕТРИЧЕСКИХ ПРЕДСТАВЛЕНИЙ МЕТОД - метод теории функций комплексного переменного, возникший из параметрического представления однолистных функций и базирующийся большей частью на Лёвнера уравнении и его обобцениях (см. [1]). Самим К. Лёвнером (К. Löwner) П. П. м. использовался на классе $S$ всех регулярных однолистных в единичном круге функций $w=f(z), f(0)=0, f^{\prime}(0)=1$, для оценки коэффициентов разложений



II



\[

w=z+c_{2} z^{2}+\ldots+c_{n} z^{n}+\ldots

\]



\[

z=f^{-1}(w)=w+b_{2} w^{2}+\ldots+b_{n} w^{n}+\ldots

\]



(см. Бибербаха гипотеза). Затем П. п. м. систематически применял Г. М. Голузин при решении проблем искажения, вращения, взаимного роста и других геометрич. характеристик отображения $w=f(z)$, связанных со значениями $f\left(z_{0}\right)$ и $f^{\prime}\left(z_{0}\right)$ при фиксированном $z_{0},\left|z_{0}\right|<1$.



П. П. м. связан с теорией оптимальных процессов. Эта связь базируется на том факте, что аналитически все упомянутые выше задачи формулируются как экстремальные задачи для управляемой системы обыкновенных дифференциальных уравнений, получаемой из уравнения Јёвнера. Использование принципа максимума Понтрягина (см. Понтрягина принцип максимума) и изучение свойств функции Понтрягина по- зволяют изучить ряд новых задач на классе $S$ и его подклассах вплоть до их полвого решения либо получить результаты, сравнимые (напр., в проблеме Бибербаха) с результатами, найденными другими методами (см. [1] п. 74).



Лuт.: [1] А л е к с а н д р о в И. А., Параметрические продолжения в теории однолистных функций, М., 1976.



ПАРАМЕТРИЧЕСКОГО РЕЗОНАНСА МАТЕМАТИЧЕСКАЯ ТЕОРИЯ - раздел теории обыкновенных дифференциальных уравнений, изучающий явление параметрич. резонанса.



Пусть $S$ - нек-рая динамич. система, способная совершать лишь колебательные движения и описываемая гамильтоновой системой линейной (невозмущенным уравнением)



\[

J \dot{x}=H_{0} x \quad\left(J=\left\|\begin{array}{cc}

0 & -I_{k} \\

I_{k} & 0

\end{array}\right\|, H_{0}^{*}=H_{0}\right)

\]



с. постоянным действительным гамильтонианом $H_{0}$. Таким образом, $(2 k \times 2 k)$-матрица $J^{-1} H_{0}$ приводится к диагональному виду с чисто мнимыми элементами



\[

i \omega_{v} \quad\left(v=\mp 1, \ldots, \mp k, \omega_{-v}=-\omega_{v}\right),

\]



$\left|\omega_{v}\right|$ - собственные частоты системы. Пусть нек-рые параметры системы $S$ начинают периодически изменяться с частотой $\vartheta>0$ и малыми амплитудами, значения к-рых определяются малым параметром $\varepsilon>0$. Если возбуждения не выводят из класса линейных гамильтоновых систем, то движение системы $S$ будет описываться возмущенным уравнением



\[

J \dot{x}=\left[H_{0}+\varepsilon H_{1}(\mathfrak{\vartheta} t)+\varepsilon^{2} H_{2}(\mathfrak{\vartheta} t)+\ldots\right] x,

\]



где $H_{j}(s+2 \pi)=H_{j}(s)=H_{j}(s)^{*}, j=1,2, \ldots$, суть кусочно непрерывные, интегрируемые в $(0,2 \pi)(2 k \times 2 k)$ матрицы-функции, и ряд в правой части (1) сходится при $\varepsilon<r_{0}$, где $r_{0}$ не зависит от $t$.



Возникновение неограниченно возрастающих колебаний системы $S$ при сколь угодно малом периодич. возмущении нек-рых ее параметров наз. п а р а м е три ч е с ки м резона н с ом. Параметрич. резонанс имеет две существенные особенности: 1) спектр частот, при к-рых возникают неограниченно возрастающие колебания, не является точечным, а состоит из совокупности малых интервалов, длины к-рых зависят от амплитуды возмущений (т. е. от $\varepsilon$ ) и к-рые стягиваются в точку при $\varepsilon \rightarrow 0$; значения частот, к к-рым стягиваются эти интервалы, наз. к р и т и ч е с к и м и; 2) колебания нарастают не по степенному, а по экспоненциальному закону. Этим параметрич. резонанс намного «опаснее» (или «полезнее», в зависимости от задачи) обычного резонанса.



Пусть $i \omega_{1}, \ldots, i \omega_{k}$ - собственные значения 1-го рода, перенумерованные так, что $\omega_{1} \leqslant \omega_{2} \leqslant \ldots \leqslant \omega_{k}$. Тогда критическими могут быть лишь частоты вида



\[

\vartheta_{j h}^{(m)}=\frac{1}{m}\left|\omega_{j}+\omega_{h}\right|, j, h=1, \ldots, k ; m=1,2, \ldots

\]



Пусть собственные векторы $f_{v}$ матрицы $J^{-1} H_{0}$, для к-рых $J^{-1} H_{0} f_{v}=i \omega_{v} f_{v}, v=\mp 1, \ldots, \mp k$, нормированы так, что



\[

i\left(J f_{\nu}, f_{\mu}\right)=\delta_{v \mu} \operatorname{sign} v, \quad v, \mu=\mp 1, \ldots, \mp k,

\]



где $\delta_{v \mu}$ - символ Кронекера, а



\[

H_{1}(\vartheta t) \sim \sum e^{i l \vartheta t} H_{1}^{(l)} .

\]



Тогда области неустойчивости в первом приближении по $\varepsilon$ определяются неравенствами



\[

\mathfrak{\vartheta}_{j h}^{(m)}+\mu_{1} \varepsilon+\ldots<\vartheta<\vartheta_{j h}^{(m)}+\mu_{2} \varepsilon+\ldots,

\]





\end{document}